\section{Setup}

Prediction-powered inference (PPI) considers a setting in which
gold-standard labels are expensive, while predictions from a machine
learning model are available at much larger scale.

Let
\[
(X,Y) \sim P,
\]
where $X \in \mathcal{X}$ denotes the covariates and $Y \in \mathcal{Y}$
denotes the gold-standard outcome.

The goal is to conduct valid statistical inference on a population
parameter
\[
\theta^* = \theta(P),
\]
rather than to improve predictive accuracy itself.

Typical examples include:
\[
\E[Y],
\qquad
\text{a quantile of } Y,
\qquad
\text{a linear regression coefficient},
\qquad
\text{a logistic regression coefficient}.
\]

We observe two samples.

\subsection{Labeled sample}

The labeled sample consists of
\[
(X_i,Y_i), \qquad i=1,\ldots,n,
\]
where both $X_i$ and the gold-standard outcome $Y_i$ are observed.

Denote
\[
\mathcal{D}_{L}
=
\{(X_i,Y_i)\}_{i=1}^n.
\]

\subsection{Unlabeled sample}

The unlabeled sample consists conceptually of
\[
(\widetilde X_i,\widetilde Y_i),
\qquad i=1,\ldots,N,
\]
but only $\widetilde X_i$ is observed.

Typically,
\[
N \gg n.
\]

A prediction rule
\[
f:\mathcal{X}\rightarrow\mathcal{Y}
\]
can be applied to both samples, giving
\[
f(X_i)
\]
for the labeled observations and
\[
f(\widetilde X_i)
\]
for the unlabeled observations.

The central idea of PPI is
\[
\boxed{
\text{large-scale predictions estimate the main signal}
+
\text{small-scale labels estimate the prediction-induced error}.
}
\]

Thus, predictions are treated as auxiliary information rather than as
gold-standard outcomes.

% ============================================================
\section{Base Assumptions}

The basic formulation of PPI relies on the following setup.

\subsection{Common population}

The labeled and unlabeled observations arise from the same target
population:
\[
(X_i,Y_i) \overset{\mathrm{iid}}{\sim} P,
\]
and conceptually
\[
(\widetilde X_i,\widetilde Y_i)
\overset{\mathrm{iid}}{\sim} P.
\]

Thus, in the base setting,
\[
P_L=P_U=P.
\]

In particular, there is no distribution shift between the labeled and
unlabeled samples.

\subsection{Independence of labeled and unlabeled samples}

The two samples are mutually independent:
\[
\mathcal{D}_L
\perp\!\!\!\perp
\mathcal{D}_U.
\]

\subsection{Fixed prediction rule}

The prediction rule $f$ is assumed to be independent of the inference
samples. For example, $f$ may be trained on a separate external dataset.

Thus, conditional on the trained predictor, $f$ can be treated as fixed
during inference.

\subsection{No correctness assumption on the prediction model}

PPI does not require
\[
f(X)=\E[Y\mid X],
\]
nor does it require
\[
\E[f(X)-Y]=0.
\]

The prediction model may be biased, misspecified, or otherwise imperfect.

The key distinction is:
\[
\boxed{
\text{prediction quality primarily affects efficiency, not validity}.
}
\]

Validity is obtained by explicitly correcting the error induced by
replacing $Y$ with $f(X)$.

% ============================================================
\section{Mean Estimation}

Consider first the simplest target:
\[
\theta^*=\E[Y].
\]

\subsection{Classical estimator}

Using only the labeled sample,
\[
\widehat{\theta}^{\class}
=
\frac{1}{n}\sum_{i=1}^n Y_i.
\]

Its variance is
\[
\Var(\widehat{\theta}^{\class})
=
\frac{\Var(Y)}{n}.
\]

This estimator is valid but may be inefficient when $n$ is small.

\subsection{Naive imputation}

A natural but generally invalid estimator is
\[
\widetilde{\theta}^{\,f}
=
\frac{1}{N}\sum_{i=1}^N f(\widetilde X_i).
\]

Although this quantity may have very small sampling variance when $N$ is
large,
\[
\E[\widetilde{\theta}^{\,f}]
=
\E[f(X)],
\]
which need not equal $\E[Y]$.

Thus, a large unlabeled sample cannot eliminate systematic prediction
bias.

\subsection{Rectifier}

Define the mean rectifier as
\[
\Delta
=
\E[f(X)-Y].
\]

Then
\[
\E[Y]
=
\E[f(X)]-\Delta.
\]

The labeled sample provides an estimator of the rectifier:
\[
\widehat{\Delta}
=
\frac{1}{n}
\sum_{i=1}^n
\{f(X_i)-Y_i\}.
\]

Hence the prediction-powered estimator is
\[
\boxed{
\widehat{\theta}^{\PP}
=
\frac{1}{N}\sum_{i=1}^N f(\widetilde X_i)
-
\frac{1}{n}\sum_{i=1}^n
\{f(X_i)-Y_i\}.
}
\]

Equivalently,
\[
\widehat{\theta}^{\PP}
=
\overline f_U-\overline f_L+\overline Y_L,
\]
or
\[
\boxed{
\widehat{\theta}^{\PP}
=
\overline Y_L
+
(\overline f_U-\overline f_L).
}
\]

This second representation shows that PPI can be viewed as a classical
estimator plus an auxiliary-information adjustment.

\subsection{Validity}

Since
\[
\E[\overline f_U]=\E[f(X)]
\]
and
\[
\E[\widehat{\Delta}]
=
\E[f(X)-Y],
\]
we have
\[
\E[\widehat{\theta}^{\PP}]
=
\E[f(X)]
-
\E[f(X)-Y]
=
\E[Y].
\]

Therefore,
\[
\boxed{
\E[\widehat{\theta}^{\PP}]=\theta^*
}
\]
even when the predictor is systematically biased.

\subsection{Efficiency}

Because the labeled and unlabeled samples are independent,
\[
\boxed{
\Var(\widehat{\theta}^{\PP})
=
\frac{\Var(f(X))}{N}
+
\frac{\Var(f(X)-Y)}{n}.
}
\]

By comparison,
\[
\Var(\widehat{\theta}^{\class})
=
\frac{\Var(Y)}{n}.
\]

When $N\gg n$,
\[
\frac{\Var(f(X))}{N}
\]
is often negligible, and hence
\[
\Var(\widehat{\theta}^{\PP})
\approx
\frac{\Var(f(X)-Y)}{n}.
\]

Thus, when $f(X)$ captures substantial variation in $Y$,
\[
\Var(f(X)-Y)
<
\Var(Y),
\]
and the prediction-powered estimator is more efficient.

The essential variance-reduction mechanism is therefore
\[
\boxed{
\Var(Y)
\quad\longrightarrow\quad
\Var(f(X)-Y)
}
\]
for the component that must be estimated using the scarce labeled sample.

A corresponding asymptotic confidence interval is
\[
\widehat{\theta}^{\PP}
\pm
z_{1-\alpha/2}
\sqrt{
\frac{\widehat{\sigma}_{f-Y}^2}{n}
+
\frac{\widehat{\sigma}_{f}^2}{N}
}.
\]

% ============================================================
\section{General Estimation}

The mean example can be generalized to a broad class of convex
$M$-estimation problems.

\subsection{Population optimization problem}

Suppose the target parameter is defined by
\[
\boxed{
\theta^*
\in
\arg\min_{\theta\in\Theta}
\E[\ell_\theta(X,Y)],
}
\]
where $\ell_\theta(x,y)$ is convex in $\theta$.

Examples include squared loss, logistic loss, and quantile loss.

Let
\[
g_\theta(x,y)
\in
\partial_\theta \ell_\theta(x,y)
\]
denote a gradient or subgradient.

Under a suitable nondegeneracy condition, the target parameter satisfies
the population estimating equation
\[
\boxed{
\E[g_{\theta^*}(X,Y)]=0.
}
\]

Define
\[
\Psi(\theta)
=
\E[g_\theta(X,Y)].
\]

Then
\[
\Psi(\theta^*)=0.
\]

\subsection{Prediction-based estimating equation}

Replacing $Y$ by the prediction $f(X)$ gives
\[
g_\theta(X,f(X)).
\]

Define the prediction-based population estimating function
\[
\boxed{
g_\theta^f
=
\E[g_\theta(X,f(X))].
}
\]

In general,
\[
g_\theta^f
\neq
\E[g_\theta(X,Y)],
\]
so solving
\[
g_\theta^f=0
\]
may lead to invalid inference.

\subsection{General rectifier}

Define the rectifier
\[
\boxed{
\Delta_\theta
=
\E\left[
g_\theta(X,Y)
-
g_\theta(X,f(X))
\right].
}
\]

Then the true population estimating equation admits the exact
decomposition
\[
\boxed{
\E[g_\theta(X,Y)]
=
g_\theta^f
+
\Delta_\theta.
}
\]

At the target parameter,
\[
\boxed{
g_{\theta^*}^f
+
\Delta_{\theta^*}
=
0.
}
\]

This is the central identity underlying general prediction-powered
inference.

The key point is that PPI does not generally correct the parameter
$\theta$ directly. Rather, it corrects the estimating equation:
\[
\boxed{
\text{true estimating equation}
=
\text{prediction-based estimating equation}
+
\text{rectifier}.
}
\]

\subsection{Sample estimators}

The prediction component can be estimated using the large unlabeled
sample:
\[
\boxed{
\widehat g_\theta^f
=
\frac{1}{N}
\sum_{i=1}^N
g_\theta
\bigl(
\widetilde X_i,
f(\widetilde X_i)
\bigr).
}
\]

The rectifier is estimated using the labeled sample:
\[
\boxed{
\widehat{\Delta}_\theta
=
\frac{1}{n}
\sum_{i=1}^n
\left[
g_\theta(X_i,Y_i)
-
g_\theta(X_i,f(X_i))
\right].
}
\]

Thus the prediction-powered estimating function is
\[
\boxed{
\widehat\Psi^{\PP}(\theta)
=
\widehat g_\theta^f
+
\widehat{\Delta}_\theta.
}
\]

The target parameter should satisfy
\[
\widehat\Psi^{\PP}(\theta^*)
\approx 0.
\]

\subsection{Confidence-set construction}

For each candidate $\theta$, one may construct a confidence region
$R_\delta(\theta)$ for the rectifier:
\[
\Pr\left(
\Delta_\theta
\in
R_\delta(\theta)
\right)
\ge
1-\delta,
\]
and a confidence region $T_{\alpha-\delta}(\theta)$ for the prediction
component:
\[
\Pr\left(
g_\theta^f
\in
T_{\alpha-\delta}(\theta)
\right)
\ge
1-(\alpha-\delta).
\]

Since
\[
g_{\theta^*}^f+\Delta_{\theta^*}=0,
\]
a prediction-powered confidence set is
\[
\boxed{
C_\alpha^{\PP}
=
\left\{
\theta:
0
\in
R_\delta(\theta)
+
T_{\alpha-\delta}(\theta)
\right\},
}
\]
where $+$ denotes the Minkowski sum.

Under the stated validity conditions,
\[
\boxed{
\Pr\left(
\theta^*
\in
C_\alpha^{\PP}
\right)
\ge
1-\alpha.
}
\]

In practical implementations, the two independent components are often
combined directly through a central limit theorem:
\[
\widehat\Psi^{\PP}(\theta)
=
\widehat g_\theta^f
+
\widehat\Delta_\theta,
\]
with approximate covariance
\[
\Var\left\{
\widehat\Psi^{\PP}(\theta)
\right\}
\approx
\frac{\Sigma_\Delta(\theta)}{n}
+
\frac{\Sigma_g(\theta)}{N}.
\]

The resulting confidence set can then be viewed as an inversion of the
corrected estimating equation:
\[
\boxed{
\text{retain }\theta
\text{ whenever }
0
\text{ is statistically compatible with }
\widehat\Psi^{\PP}(\theta).
}
\]

% ============================================================
\section{Examples and Algorithms}

\subsection{Mean estimation}

For mean estimation, take
\[
\ell_\theta(y)
=
\frac{1}{2}(y-\theta)^2.
\]

Then
\[
g_\theta(y)
=
\theta-y.
\]

The prediction-based estimating function is
\[
g_\theta^f
=
\E[\theta-f(X)]
=
\theta-\E[f(X)].
\]

The rectifier is
\[
\begin{aligned}
\Delta_\theta
&=
\E[
(\theta-Y)-(\theta-f(X))
]
\\
&=
\E[f(X)-Y].
\end{aligned}
\]

Thus $\Delta_\theta$ does not depend on $\theta$, and the corrected
estimating equation is
\[
\theta-\E[f(X)]
+
\E[f(X)-Y]
=
0.
\]

Solving explicitly gives
\[
\theta^*
=
\E[f(X)]
-
\E[f(X)-Y],
\]
leading to
\[
\widehat\theta^{\PP}
=
\overline f_U
-
\overline{(f-Y)}_L.
\]

Thus mean estimation is a case in which the general estimating equation
has a closed-form solution.

% ------------------------------------------------------------
\subsection{Quantile estimation}

Let $\theta^*$ denote the $q$-quantile of $Y$. It can be represented as a
minimizer of the quantile loss.

A suitable subgradient is
\[
\boxed{
g_\theta(y)
=
\mathbf{1}\{y\leq\theta\}-q.
}
\]

Hence
\[
\E[g_\theta(Y)]
=
F_Y(\theta)-q,
\]
and the target satisfies
\[
F_Y(\theta^*)=q.
\]

The prediction component is
\[
g_\theta^f
=
F_f(\theta)-q,
\]
where
\[
F_f(\theta)
=
\Pr(f(X)\leq\theta).
\]

The rectifier is
\[
\begin{aligned}
\Delta_\theta
&=
\E\left[
\mathbf{1}\{Y\leq\theta\}
-
\mathbf{1}\{f(X)\leq\theta\}
\right]
\\
&=
F_Y(\theta)-F_f(\theta).
\end{aligned}
\]

Therefore,
\[
g_\theta^f+\Delta_\theta
=
F_Y(\theta)-q.
\]

The sample estimators are
\[
\widehat F_f(\theta)
=
\frac{1}{N}
\sum_{i=1}^N
\mathbf{1}
\{f(\widetilde X_i)\leq\theta\},
\]
and
\[
\widehat\Delta_\theta
=
\frac{1}{n}
\sum_{i=1}^n
\left[
\mathbf{1}\{Y_i\leq\theta\}
-
\mathbf{1}\{f(X_i)\leq\theta\}
\right].
\]

Thus
\[
\widehat\Psi^{\PP}(\theta)
=
\widehat F_f(\theta)
+
\widehat\Delta_\theta
-
q.
\]

Unlike the mean case, $\theta$ enters through indicator functions, so
there is generally no closed-form solution. Instead, inference is
performed by evaluating candidate values of $\theta$ and retaining those
for which
\[
\widehat\Psi^{\PP}(\theta)
\]
is statistically compatible with zero.

This example illustrates why general PPI is naturally formulated as
correction of the estimating equation rather than direct correction of
the parameter.

% ------------------------------------------------------------
\subsection{Logistic regression}

Suppose $Y\in\{0,1\}$ and define the population logistic regression
coefficient by
\[
\theta^*
=
\arg\min_\theta
\E
\left[
-YX^\top\theta
+
\log(1+\exp(X^\top\theta))
\right].
\]

Define
\[
\mu_\theta(x)
=
\frac{1}{1+\exp(-x^\top\theta)}.
\]

The gradient is
\[
\boxed{
g_\theta(x,y)
=
x\{\mu_\theta(x)-y\}.
}
\]

Hence the true population score equation is
\[
\E[
X\{\mu_{\theta^*}(X)-Y\}
]
=
0.
\]

The prediction-based component is
\[
\boxed{
g_\theta^f
=
\E[
X\{\mu_\theta(X)-f(X)\}
].
}
\]

The rectifier is
\[
\begin{aligned}
\Delta_\theta
&=
\E\left[
X\{\mu_\theta(X)-Y\}
-
X\{\mu_\theta(X)-f(X)\}
\right]
\\
&=
\E[X\{f(X)-Y\}].
\end{aligned}
\]

Thus,
\[
\boxed{
\Delta_\theta
\equiv
\Delta
=
\E[X(f(X)-Y)],
}
\]
which is independent of $\theta$.

The sample estimators are
\[
\widehat g_\theta^f
=
\frac{1}{N}
\sum_{i=1}^N
\widetilde X_i
\left[
\mu_\theta(\widetilde X_i)
-
f(\widetilde X_i)
\right],
\]
and
\[
\widehat\Delta
=
\frac{1}{n}
\sum_{i=1}^n
X_i
\{f(X_i)-Y_i\}.
\]

The corrected score is
\[
\boxed{
\widehat\Psi^{\PP}(\theta)
=
\widehat g_\theta^f
+
\widehat\Delta.
}
\]

Because $\mu_\theta(x)$ is nonlinear in $\theta$, this estimating
equation does not generally admit a closed-form solution. Confidence
sets are therefore obtained through test inversion over candidate
values of $\theta$.

For a $d$-dimensional parameter, inference can be based on the
coordinate-wise conditions
\[
\left|
\widehat\Psi_j^{\PP}(\theta)
\right|
\leq
w_{\alpha,j}(\theta),
\qquad j=1,\ldots,d,
\]
with multiplicity adjustment as needed.

% ------------------------------------------------------------
\subsection{Linear regression}

Define the population least-squares coefficient by
\[
\theta^*
=
\arg\min_\theta
\E[(Y-X^\top\theta)^2].
\]

The population normal equation is
\[
\E[
X(X^\top\theta^*-Y)
]
=
0.
\]

Let
\[
A=\E[XX^\top],
\]
and assume $A$ is nonsingular. Then
\[
\theta^*
=
A^{-1}\E[XY].
\]

A gradient is
\[
g_\theta(X,Y)
=
X(X^\top\theta-Y).
\]

The prediction-based component is
\[
g_\theta^f
=
\E[
X(X^\top\theta-f(X))
].
\]

The rectifier is
\[
\begin{aligned}
\Delta_\theta
&=
\E\left[
X(X^\top\theta-Y)
-
X(X^\top\theta-f(X))
\right]
\\
&=
\E[X(f(X)-Y)].
\end{aligned}
\]

Again, the rectifier does not depend on $\theta$.

The corrected population equation is
\[
A\theta
-
\E[Xf(X)]
+
\E[X(f(X)-Y)]
=
0.
\]

Thus
\[
\boxed{
\theta^*
=
A^{-1}\E[Xf(X)]
-
A^{-1}\E[X(f(X)-Y)].
}
\]

Define
\[
\theta^f
=
A^{-1}\E[Xf(X)]
\]
and
\[
\Delta
=
A^{-1}\E[X(f(X)-Y)].
\]

Then
\[
\boxed{
\theta^*
=
\theta^f-\Delta.
}
\]

This yields a particularly simple sample implementation:
\[
\boxed{
\widehat\theta^{\PP}
=
\widetilde X^\dagger f(\widetilde X)
-
X^\dagger\{f(X)-Y\},
}
\]
where $\dagger$ denotes the Moore--Penrose pseudoinverse.

Thus, linear-regression PPI can be interpreted as
\[
\boxed{
\text{regression coefficient of predictions}
-
\text{regression coefficient of prediction errors}.
}
\]

Because the OLS estimating equation is linear in $\theta$, a
closed-form prediction-powered estimator is available.

Its asymptotic covariance has the generic form
\[
\Var(\widehat\theta^{\PP})
\approx
\frac{V_\Delta}{n}
+
\frac{V_f}{N},
\]
where $V_\Delta$ and $V_f$ are the appropriate sandwich covariance
matrices for the rectifier and prediction-based regressions,
respectively.

% ------------------------------------------------------------
\subsection{Generic convex estimation}

For a general convex loss,
\[
\theta^*
\in
\arg\min_\theta
\E[\ell_\theta(X,Y)],
\]
let
\[
g_\theta(x,y)
\in
\partial_\theta\ell_\theta(x,y).
\]

For every candidate $\theta$, calculate
\[
\boxed{
\widehat\Delta_\theta
=
\frac{1}{n}
\sum_{i=1}^n
\left[
g_\theta(X_i,Y_i)
-
g_\theta(X_i,f(X_i))
\right]
}
\]
and
\[
\boxed{
\widehat g_\theta^f
=
\frac{1}{N}
\sum_{i=1}^N
g_\theta(
\widetilde X_i,
f(\widetilde X_i)
).
}
\]

The corrected estimating function is
\[
\boxed{
\widehat\Psi^{\PP}(\theta)
=
\widehat g_\theta^f
+
\widehat\Delta_\theta.
}
\]

Its asymptotic covariance is obtained by combining the two independent
sampling components:
\[
\boxed{
\widehat\Sigma_{\PP}(\theta)
=
\frac{\widehat\Sigma_\Delta(\theta)}{n}
+
\frac{\widehat\Sigma_g(\theta)}{N}.
}
\]

Inference proceeds by retaining candidate values of $\theta$ for which
zero is statistically compatible with the corrected estimating
function:
\[
\boxed{
C_\alpha^{\PP}
=
\left\{
\theta:
0
\text{ lies in a confidence region for }
\widehat\Psi^{\PP}(\theta)
\right\}.
}
\]

Thus, the specific PPI algorithms differ primarily in two structural
properties:

\begin{enumerate}[label=(\arabic*)]
    \item whether the rectifier $\Delta_\theta$ depends on $\theta$;
    \item whether the corrected estimating equation can be solved
    explicitly for $\theta$.
\end{enumerate}

These distinctions can be summarized as follows.

\begin{center}
\begin{tabular}{llll}
\toprule
Problem
& Rectifier
& Depends on $\theta$?
& Implementation
\\
\midrule
Mean
& $\E[f(X)-Y]$
& No
& Closed form
\\

Quantile
&
$\E[
\mathbf{1}\{Y\leq\theta\}
-
\mathbf{1}\{f(X)\leq\theta\}
]$
& Yes
& Inversion
\\

Logistic regression
&
$\E[X(f(X)-Y)]$
& No
& Inversion
\\

Linear regression
&
$\E[X(f(X)-Y)]$
& No
& Closed form
\\

General convex estimation
&
$\E[
g_\theta(X,Y)
-
g_\theta(X,f(X))
]$
& Generally yes
& Inversion
\\
\bottomrule
\end{tabular}
\end{center}